\begin{frame}[t]{\ev{\tagM\,\tagB\,\tagP}The interface the runtime owes the loop}
\vspace{0pt}
{\small
\begin{tabularx}{\textwidth}{@{}>{\ttfamily}L{7.6cm} Y@{}}
\toprule
\normalfont Call & Status\\
\midrule
checkpoint(cell) $\to$ snap & \tagM\ sandbox-platform named snapshots, memory unverified\\
fork(snap, lease, seed=copy|reseed, egress) & \tagP\\
evaluate(artifact, verifier) $\to$ receipt & \tagP\ runs outside the child\\
select(candidates, fresh test) $\to$ one & \tagP\ selection without pooling\\
reduce(receipts) $\to$ $(\hat\theta, N_{\rm eff}$ bound$)$ | abstain & \tagB\ merge + evidence check,\ \tagP\ $N_{\rm eff}$ bound, abstain\\
promote(epoch, receipt) & \tagB\ sha256-bound gate,\ \tagP\ epoch fence\\
\bottomrule
\end{tabularx}\par}
\vspace{.08cm}
{\small The epoch fence stops a stale run from writing, like a Chubby sequencer or a fencing token.\par}
\claim{Select chooses and reduce pools, so they are different calls.}
\sources{Fork contract (design document), sandbox-platform named snapshots: \href{https://arxiv.org/abs/2607.09689}{arXiv:2607.09689} \S4; Burrows OSDI'06; Kleppmann (2016). Architecture: next slide.}
\note{[0:55] Together, these give the interface the runtime owes the loop. The six calls carry a stated status. Checkpoint uses named snapshots with memory unverified. Fork takes a lease, an egress policy and a seed policy, copy or reseed. Evaluate outside the child. Select one candidate on a fresh test. Reduce combines receipts into an estimate or abstains, with merge built and abstention proposed. Promote runs once, with the hash-bound gate built and the epoch fence, which blocks stale runs, proposed. Select and reduce are separate calls.}
\end{frame}
